%% Document metadata — title, authors, keywords.
%%
%% THIS FILE IS YOURS. It is deliberately template-agnostic: `make set-template`
%% rewrites main.tex but never touches this file, so switching between IEEE and
%% arXiv formats keeps your title and author list.
%%
%% Each template composes these macros into its own author block, because
%% \IEEEauthorblockN and authblk have nothing in common. Adding an author means
%% uncommenting the macros below AND adding one line to main.tex — see
%% tools/templates/README.md, "Author blocks".

\newcommand{\PaperTitle}{An Example Paper for the Paper Template}

%% Running head. Keep it short — it has to fit across the page.
\newcommand{\PaperShortTitle}{Paper Template}

%% Comma-separated. Used for \IEEEkeywords and arxiv.sty's \keywords.
\newcommand{\PaperKeywords}{reproducibility, research infrastructure, LaTeX}

%% PDF subject line — arXiv categories work well here, e.g. cs.LG, cs.AI.
\newcommand{\PaperSubject}{cs.DL}

%% Flat author list for the PDF metadata.
\newcommand{\PaperAuthorList}{Author Name}

%% --- Authors ----------------------------------------------------------------
\newcommand{\AuthorOneName}{Author Name}
\newcommand{\AuthorOneAffil}{Institution}
\newcommand{\AuthorOneEmail}{author@example.org}
\newcommand{\AuthorOneOrcid}{https://orcid.org/0000-0000-0000-0000}

% \newcommand{\AuthorTwoName}{Second Author}
% \newcommand{\AuthorTwoAffil}{Institution}
% \newcommand{\AuthorTwoEmail}{second@example.org}
% \newcommand{\AuthorTwoOrcid}{https://orcid.org/0000-0000-0000-0000}

% \newcommand{\AuthorThreeName}{Third Author}
% \newcommand{\AuthorThreeAffil}{Institution}
% \newcommand{\AuthorThreeEmail}{third@example.org}
% \newcommand{\AuthorThreeOrcid}{https://orcid.org/0000-0000-0000-0000}
